\scanpage{22}\mspage{20}
dans un cas, trois dans l'autre).

Dans le cas général, il y a lieu d'expliciter
la relation d'ordre de \unsure{la|l'} \deleted{\unsure{catégorie}} $\underline{\mathrm{Hom}}^{!}(\unsure{\mathscr{L}|\mathcal{L}}_k,\mathscr{L}_l)$
($!$ signifie : compatible aux source et but), en termes
des \deleted{$\unsure{E_\varphi}$} ensembles représentatifs $E_\varphi,E_\psi$.
Je ne vais pas expliciter cette relation d'ordre, mais
seulement une propriété de celle-ci qui nous sera
suffisante. \deleted{\unsure{Notons de}} Soit $\varphi_E$ le foncteur défini
par une partie $E\subset[1,2k]$. Tout d'abord :

\uline{Lemme \unsure{1|2|3}} (immédiat) : \inserted{Soient $E,E'\subset[1,2k]$ admissibles, et} soit $j\in[1,2k]$ tel que
$E\cap[1,j]=E'\cap[1,j]$. Alors $\varphi_E$ et $\varphi_{E'}$ coïncident sur \unsure{les}
le \deleted{sommet} \inserted{but} de $u_j$ ($j$.\unsure{ième} flèche de $\mathscr{L}_k$) et sur
tous les objets de $\mathscr{L}_k$ \deleted{\unsure{de rang}} antérieurs à celle-ci
(au sens de la numérotation des objets, et non de la
relation d'ordre de la catégorie $\unsure{\mathscr{L}}_k$ bien sûr). Dualement,
si $E\cap[j,2k]=E'\cap[j,2k]$, alors $\varphi_E$ et $\varphi_{E'}$ coïncident
sur la source de $u_j$, et sur tous les objets de $\mathscr{L}_k$ postérieurs
à celui-ci.

Considérons maintenant une \deleted{composante connexe} \inserted{partie connexe}
$I$ de $E$, \deleted{et soit $I'$} (pas forcément une composante
connexe), \unsure{soit} $I=[r,s]\subset[1,2k]$ ($r$ le plus petit élément de $I$,
$s$ son \deleted{élément} plus grand élément.) et supposons que

a) $\operatorname{card}I$ est pair (i.e. $s-r$ est pair)

b) $r\ne1$, et $r-1\notin E$ (i.e. $r-1\in[1,2k]\setminus E$)

et soit \deleted{$\unsure{E'=}$}
\[
E'=(E\setminus\{s\})\cup\{r-1\}=(E\setminus I)\cup\overbrace{[r-1,s-1]}^{I'}
\]
\unsure{dérivant|déduit} de $I$ en ``translatant'' $I$ sur la gauche d'un cran.
% archive sheet number: 21
